% BHU PYQ -- Professional Exam Template
\documentclass[11pt,a4paper]{article}

\usepackage[a4paper,top=2.5cm,bottom=2.5cm,left=2.8cm,right=2.8cm,headheight=14pt]{geometry}
\usepackage{amsmath,amssymb}
\usepackage[T1]{fontenc}
\usepackage[utf8]{inputenc}
\usepackage{lmodern}
\usepackage{microtype}
\usepackage{fancyhdr}
\usepackage{enumitem}
\usepackage{listings}
\usepackage{hyperref}
\usepackage{lastpage}
\usepackage{parskip}

\hypersetup{colorlinks=true,urlcolor=black,linkcolor=black,
  pdftitle={BVCA-303: Programming In Linux -- BHU PYQ}}

\pagestyle{fancy}
\fancyhf{}
\renewcommand{\headrulewidth}{0.4pt}
\renewcommand{\footrulewidth}{0.4pt}
\fancyhead[L]{\small\textit{Banaras Hindu University}}
\fancyhead[R]{\small\textit{BVCA-303: Programming In Linux}}
\fancyfoot[C]{\small Page \thepage\ of \pageref{LastPage}}

\lstset{
  basicstyle=\small\ttfamily,
  frame=single,
  breaklines=true,
  columns=flexible,
  keepspaces=true
}

\newlist{parts}{enumerate}{2}
\setlist[parts,1]{label=(\alph*),leftmargin=*,itemsep=4pt,topsep=3pt}
\setlist[parts,2]{label=(\roman*),leftmargin=*,itemsep=2pt,topsep=2pt}
\newcommand{\pts}[1]{\hfill\textit{[#1]}}

\begin{document}
\begin{center}
  {\large\bfseries BANARAS HINDU UNIVERSITY}\\[2pt]
  {\normalsize B.Voc. Semester III Examination, 2023-24}\\[6pt]
  \rule{0.8\linewidth}{0.5pt}\\[6pt]
  {\large\bfseries Paper: BVCA-303 --- Programming In Linux}\\[4pt]
  {\normalsize Time: Three hours \hfill Maximum Marks: 70}
\end{center}

\rule{\linewidth}{0.4pt}

\smallskip
\noindent\textit{Instructions: Answer the questions in each section as instructed. All questions carry marks as indicated.}

\bigskip

TROMIN Ossian stettcattrskonessseossoeaseaplity
BVCA-303 : Programming in Linux

\bigskip\noindent{\large\bfseries SECTION --- A}
\rule{\linewidth}{0.4pt}\smallskip

Note: Long-Answer Questions each in 500 words approximately (Answer any 2 out of 2x12=24
4; Each question carries 12 marks.
!) Explain the architecture of Linux, emphasizing the roles of the kernel, shell, and file system.

\medskip\noindent\textbf{Question 2.} Describe the Linux file system hierarchy in detail, providing examples of directories at each level.
Explain the purpose and usage of essential file system management commands such as Is, cp, mv, and
rm. Illustrate their functionality with practical examples,

\medskip\noindent\textbf{Question 3.} Discuss the significance of Linux commands in performing various tasks efficiently. Explain how
commands like touch, grep, we, and cat contribute to file manipulation, searching, and text
processing. Provide examples to demonstrate the practical utility of these commands in everyday naan
computing tasks.

\medskip\noindent\textbf{Question 4.} Compare and contrast the advantages and disadvantages of using Linux as an operating system.
Provide real-world examples to support your arguments.

\bigskip\noindent{\large\bfseries SECTION --- B}
\rule{\linewidth}{0.4pt}\smallskip

\medskip\noindent\textbf{Question 5.} Answer any six of the following in about 200 words each : 6x6=36
| (a) Write a shell script to print a pattern of asterisks (*) in the following manner:
| *
| *
| Ke OK
\begin{parts}
  \item Develop a shell script that accepts two numbers as input and performs addition, subtraction,
  multiplication, and division on them.
  (C) Discuss the significance of pipes in Linux. Provide an example scenario where pipes are used
  effectively with multiple commands. ’ :
  \item Discuss the purpose of the mkdir command in Linux. Provide an example demonstrating its usage
  in creating directories.
  (€) Define the tee command in Linux and explain its functionality.
  \item What is the purpose of the cat command in Linux? How does it differ from other file viewing
  commands?
  \item Explain the significance of the banner command in Linux. How can it be used to display text in a
  visually appealing manner?
  \item Discuss the functionality of the grep command in Linux.
  \_(i) Describe the significance of the cal command in Linux. How does it differ from other calendar-
  related commands? Provide an example demonstrating its usage.
  P.T.0.
\end{parts}

\bigskip\noindent{\large\bfseries SECTION --- C}
\rule{\linewidth}{0.4pt}\smallskip

\medskip\noindent\textbf{Question 6.} Answer the following questions in a word or sentence each: 1x10=10
\begin{parts}
  \item Which command is used to display the current working directory in Linux? i ,
  : (a) Is (b)PWD (c)cal (d)cp
  \begin{parts}
    \item Which command is used to create a new directory in Linux?
  \end{parts}
  \item mkdir (b)touch (c)rm  (d)mv
  \begin{parts}
    \item What does the man command do in Linux?
  \end{parts}
  \item Displays the manual page of acommand (b) Moves files or directories
  \item Creates a new file (d) Counts the number of words in a file
  \begin{parts}
    \item What does the head command do in Linux?
  \end{parts}
  \item Displays the first part of a file (b) Displays the last part of a file
  \item Searches for a pattern in a file (d) Counts the number of lines in a file
  \begin{parts}
    \item Which command is used to count lines in a file in Linux?
  \end{parts}
  \item we (b) sort (c) grep (d) banner |
  \begin{parts}
    \item What is the purpose of the grep command in Linux?
  \end{parts}
  \item Displays the calendar for the current month (b) Displays the current date and time
  \item Counts the number of lines in a file (d) Searches for a pattern in a file
  \begin{parts}
    \item Which command is used to count the number of words in a file in Linux?
  \end{parts}
  \item wv (b) we (c) ps (d) tee
  \begin{parts}
    \item What is the function of the sort command in Linux?
  \end{parts}
  \item Displays the sorted content of a file (b) Moves files or directories
  \item Counts the number of lines in a file (d) Searches for a pattern in a file
  \begin{parts}
    \item Which command is used to display the last part of a file in Linux?
  \end{parts}
  \item head (b) tail (c) who (d) be
  \begin{parts}
    \item What does the mkdir command do in Linux?
  \end{parts}
  \item Removes a directory (b) Creates a new directory
  \item Moves files or directories (d) Displays the manual page of a command
  ARK
\end{parts}

\vfill
\rule{\linewidth}{0.4pt}
\begin{center}\small
\href{https://bhu-exam.vercel.app/}{Click here to visit our website}\\[4pt]\href{https://me-aryan.vercel.app/}{Click here to visit our contributor}\\[4pt]\href{https://quashan-me.vercel.app/}{Click here to visit our contributor}
\end{center}
\end{document}